% !TEX root = ./main.tex
\documentclass[a4paper]{article}
% Course language: use [english] or [german].
\usepackage[german]{ajnotes}
\title{\textbf{Betriebwirtschaftslehre — Marketing A}}
\begin{document}
    \maketitle
    % Attach the given problem sheet: drop it in this folder as sheet.pdf and
    % un-comment the next line (pdfpages is already loaded by the preamble).
    % \includepdf[pages=-]{sheet.pdf}
\section{Admin}
\begin{itemize}
    \item Übungen vor der Lektion durchlesen
    \item Bei Fragen: ergenzinger@business.uzh.ch
    \item Keine Musterlösungen
    \item \textbf{Mahalo Case Study bis 2. übung lesen}
\end{itemize}

\section{1. Übungseinheit}
4. ersten Säulen des Marketingprozesses werden anhand der Unternehmenswertschöpfung zusammen kategorisiert
5. Säule ist Seite des Kunden gegenleistung für Produkte

\subsection{SOR Modell}
Stimuli-Organismus-Response, wobei Stimuli und Reponses unmittelbar beobachtbar sind. 
Organismus oder wie der Konsument die Informationen verarbeitet und interpretiert ist von besondere Wichtigkeit für Unternehmen.
Wenn Unternehmen

\subsection{Effektivität vs Effizienz}
Effektivität beschreibt binär, ob ein Ziel erfolgreich erreicht wurde oder nicht. Effizienz ist ein quantitatives Mass an Aufwand/Resultat, oder wie produktiv das Ziel erreicht wurde ohne zu sagen, ob das Ziel überhaupt das richtige war.

\subsection{Mahalo Case Study}
Mineralwasser aus Hawaii in die Schweiz importiert. Verkaufszahlen sind sinkend, obwohl nachfrage nach Mineralwasser steigt.
\pagebreak
\subsubsection{Interviews und Forschungsmethoden}
\begin{itemize}
    \item Weil Verkaufszahlen zurück möchte man verstehen wieso
    \item Es wurden Interviews mit 3 Personen geführt, die verschiedener Kundensegmente gehören
    \item Frage: Um welche Art von den 3 Marktforschungsmethoden geht es? (Deskriptiv, Explikativ, ) 
    \begin{itemize}
        \item Explorativ, weil es darum geht eine noch unbeantwortete Frage zu lösen
    \end{itemize} 
    \item Drei Informationen 
    \begin{itemize}
        \item Demografik dient der Segmentierung
        \item Kauf gründe dient Abgrenzung in konvergente Nutzen (Identifizierung von Wettbewerber) und divergente Nutzen
        \item Werte der Kunden dient dem Verständnis des Konsumentenverhaltens
    \end{itemize}
    \item Das Mahalo-Wasser befriedigt verschiedene Bedürfnisse (Durst, Gesundheit, Status)
    \begin{itemize}
        \item Wie heisst dieses Konzept? divergente Produktnutzen
        \item Was dient dieses Konzept? Differenzierung der Kundengruppen und Kaufgründe für gezielte Marketing Instrumente auf die Zielgruppen. 
    \end{itemize}

\end{itemize}

\subsubsection{Segmentierung}
\begin{itemize}
    \item Alle 3 Kunden können im Schwartz-Wertediagramm eingeordnet werden
\end{itemize}

\subsubsection{Marktforschung und ihre Phasen}
\begin{enumerate}
    \item Problemstellung und Zieldefinition
    \begin{itemize}
        \item Unklar welche Kundentypen das Wasser kaufen
        \item Art der Forschung: Deskriptiv. Es muss erhoben werden was die Merkmale der Kunden sind, die schon identifiziert wurden mit den Interviews.
        \item Ziel: Beschreibung der Kundentypen anhand Informationen als Basis für die Segmentierung.
    \end{itemize}
    \pagebreak
    \item Entwicklung Untersuchungsplan
    \begin{itemize}
        \item Art der Datenerhebung (\textbf{Befragung} (Fragebogen Umfrage), Beobachtung, Experimente)
        \item Kommunikationsmittel (Internet, Zeitung, Plakate, Tür zu Tür)
        \item Inhalt/Gesuchte Informationen (Kaufmotive, Werte, Gewohnheiten, Demografik)
        \item Stichprobenplan:
        \begin{itemize}
            \item Welche Subjekte? Unbekannt also zufällige Teilerhebung 
            \item Grösse? X
            \item Verfahren? (Einfache Auswahl, geschichtete Auswahl, \textbf{Flächenauswahl})
        \end{itemize}
    \end{itemize}
    \item Datenerhebung und Analyse
    \begin{itemize}
        \item Gütekriterien (Obj,Rel,Val)
        \item Zusammenfassung, Visualisierung, statistische Segmentierung
    \end{itemize}
    \item Interpretation und Kommunikation der Ergebnisse
    \begin{itemize}
        \item Diagramm mit Anzahl der Antworten nach Zeit, die gebraucht wurde um zu beantworten
        \item Sehr klein, meist sehr schnell beantwortet, also low Involvement
        \item Diagramm Grund wieso gekauft wird
        \begin{itemize}
            \item Für Mahalo Wasser meist Werbung also passive Aufnahme von Produktinformationen
        \end{itemize}
        \item Kaufentscheidung nach Unterschiede und Involvement
        \begin{itemize}
            \item wenige Unterschiede zu anderen Wassermarken und Low-Involvement
            \item also habitualisiertes Kaufverhalten
            \item \textbf{Wie müsste Mahalo am Besten verkauft werden? }
            \begin{itemize}
                \item Nach Auffällige Platzierung
                \item Starke Dispersion 
                \item Stück Aktionen
                \item Wiedererkennbarkeit fördern (fokus auf visuellen Aspekte)
            \end{itemize}
        \end{itemize}
    \end{itemize}
\end{enumerate}

\subsection{Prüfungsaufgabe}
\begin{itemize}
    \item Case BeYoutiful
    \subitem C und D weil A und B sich auf konvergente Nutzenanalyse Beziehungen
\end{itemize}
\end{document}
